\documentclass[11pt]{ctexart}
\usepackage[a4paper,margin=1in]{geometry}
\usepackage{graphicx}
\usepackage{xcolor}
\usepackage{hyperref}
\definecolor{refgray}{gray}{0.45}
\title{\textbf{Market Views｜全球投行叙事汇编}\\[0.4em]{\large AI基础设施投资从GPU向存储、光互联、配电全面扩散，消费与医疗板块结构性分化加剧，宏观增长动能切换中政策预期调整}}
\author{KC桌面 自动生成}
\date{260719}
\begin{document}
\maketitle
\tableofcontents
\newpage
\section{一页摘要}
\begin{itemize}
  \item AI基础设施投资正从GPU集群扩展至存储、光互联、配电架构及自动驾驶等下游应用，各环节均出现结构性拐点，订单能见度延伸至2028年。
  \item 存储价格走势出现分歧：BofA认为三季度DRAM ASP环比+21\%超预期，非LTA部分占比60-70\%推高弹性；TrendForce预测NAND ASP将从峰值回落至4Q26的0-5\% QoQ。
  \item 中国新能源车市场折扣分化，新势力折扣收敛而传统品牌扩大，市场对定价权是否回归存在分歧。
  \item 北美卡车订单122\%增长，市场对其性质（终端需求复苏 vs 渠道补库存）存在分歧，决定周期高度判断。
  \item 美国管理式医疗龙头联合健康超预期但股价高开低走，反映市场对盈利可持续性的定价已前置；商业险利润率恢复时间线拉长至2027年后。
  \item 中国生物科技板块估值修复主要由资金轮动驱动，P/S从3.0倍回升至3.8倍，流动性占比从6\%升至9.7\%，持续性待观察。
  \item 中国7月出口量增速因台风扰动阶段性放缓，但运价上涨反映供给端紧张未缓解；地产量缩正向价跌传导，城市层级分化明显。
  \item 锂电产业链库存增速连续超过销量增速，指向主动去库与价格承压；但7月锂电池排产创新高且储能占比突破40\%，需求结构改善可能缓解过剩压力。
\end{itemize}
\section{全球投行叙事汇编}
今日纳入 27 篇投行/券商研究，按 4 个市场、资产与产业主题系统整理。
\newpage
\subsection{AI基础设施：从芯片到配电的全面重构}
\textbf{AI基础设施投资正从GPU集群扩展至存储、光互联、配电架构及自动驾驶等下游应用，各环节均出现结构性拐点：存储受益于HBM与DDR5代际升级，光模块进入1.6T/3.2T高速迭代，数据中心配电向800VDC转型，中国AI芯片产能约束下国产化加速，自动驾驶软件商Momenta凭借跨品牌合作占据独立NOA市场64.5\%份额。}
\subsubsection*{主流共识}
\begin{itemize}
  \item AI需求持续拉动存储、光模块、电力设备等基础设施投资，订单能见度延伸至2028年
  \item 数据中心功率密度指数级增长（10kW→1MW+），传统配电架构面临根本性变革
  \item 中国AI芯片市场供需缺口至少持续至2027年，产能而非需求是核心约束
\end{itemize}
\subsubsection*{分歧与条件差异}
\begin{itemize}
  \item 存储价格走势：BofA认为三季度DRAM ASP环比+21\%超预期，非LTA部分占比60-70\%推高弹性；但TrendForce预测NAND ASP将从峰值85-90\% QoQ回落至4Q26的0-5\% QoQ
  \item 韩国电池企业关注点：JPM指出市场焦点已从订单规模转向利润拐点与闲置资产盘活，LG新能源下半年出货量弹性（10.6→19.4GWh）是关键
  \item 澜起科技增长驱动：GS强调DDR5 Gen-3/Gen-4代际升级是净利润超预期48\%的核心，而非规模效应；MS则更关注AI推理与PCIe Retimer在中国AI服务器市场的独特地位
  \item 自动驾驶软件估值：UBS认为Momenta 64.5\%市占率（独立NOA）支撑高PS，但竞争（华为、地平线）与车企压价风险不可忽视
\end{itemize}
\subsubsection*{机构观点}
\paragraph{BofA} 三季度DRAM平均售价环比增长21\%，高于市场预期，主要因非长期协议部分（占60-70\%）价格弹性被低估。服务器DRAM由LPDDR5领涨20-30\%，DDR5/DDR4现货连续八周上涨，OEM为旺季备货加大采购。台积电与ASML指引为存储需求提供额外支撑，HBM/DRAM资本开支传导至光刻环节。
{\small\color{refgray}数据：3Q26 DRAM ASP环比+21\%\par}
{\small\color{refgray}数据：非LTA DRAM销售占比60-70\%\par}
{\small\color{refgray}数据：服务器DRAM环比+20-30\%\par}
{\small\color{refgray}数据：DDR5现货价格指向约20\% QoQ涨幅\par}
{\small\color{refgray}边际变化：市场此前低估非长期协议价格弹性，OEM从观望转向备货，现货连续八周上涨信号意义增强\par}
\paragraph{Citi} ASML订单能见度延伸至2028年，2027年Low-NA EUV产能接近售罄，2028年实质性订单已流入。存储客户（HBM/DRAM）是增量主力，驱动2026-2028年系统销售预测上调16-45\%，EPS上调25-52\%。同时，AI数据中心配电架构正从传统交流转向800VDC，预计2030年覆盖79\%新增容量，固态变压器需求三年增长40倍（2027年901 MVA→2030年37,152 MVA）。光互联市场2025-2030年CAGR 31.6\%，中际旭创以21.2\%市占率领先，1.6T产品2025年起步占4\%，2030年将占45\%。
{\small\color{refgray}数据：ASML 2027年Low-NA EUV产能接近售罄\par}
{\small\color{refgray}数据：2026-2028年系统销售预测上调16-45\%\par}
{\small\color{refgray}数据：EPS上调25-52\%\par}
{\small\color{refgray}数据：800VDC 2030年覆盖79\%新增数据中心容量\par}
{\small\color{refgray}边际变化：ASML订单从季度指引拉长至三年规划，存储资本开支成为增量主力；数据中心配电从渐进升级变为结构性拐点，800VDC成为新标准\par}
\paragraph{DB} 中国AI芯片市场2025年出货约400万颗，2026年增至500万颗，供需缺口至少持续至2027年。制造产能（中芯国际、华虹）是核心约束，国产化率从40\%向50\%以上爬升。竞争格局分层：华为独居第一梯队，供应链能力成为分水岭。推理与训练装机比例约10:1，国产芯片在预训练环节仍有明显短板。
{\small\color{refgray}数据：2025年国产AI芯片出货约400万颗\par}
{\small\color{refgray}数据：2026年预计增至500万颗\par}
{\small\color{refgray}数据：国产化率从40\%提升至50\%以上\par}
{\small\color{refgray}数据：推理与训练装机比例约10:1\par}
{\small\color{refgray}边际变化：竞争焦点从硬件设计转向供应链交付能力，华为凭借产能优先级维持领先\par}
\paragraph{GS} 澜起科技2Q26净利润指引中值10-12亿元，同比+61-77\%，超出预期约48\%，核心驱动来自DDR5 Gen-3/Gen-4接口芯片占比提升及PCIe Retimer强劲增长。营收19亿元同比+34\%，毛利率从67.3\%上调至67.6\%，净利率从46.9\%大幅上调至50.0\%。Gen-5已开始出货，DDR6 Gen-1在研，MRCD/MDB Gen-2预计2027年加速采用。
{\small\color{refgray}数据：2Q26净利润指引中值10-12亿元，同比+61-77\%\par}
{\small\color{refgray}数据：超出GS预期约48\%\par}
{\small\color{refgray}数据：营收19亿元，同比+34\%\par}
{\small\color{refgray}数据：毛利率上调至67.6\%，净利率上调至50.0\%\par}
{\small\color{refgray}边际变化：增长逻辑从DDR5渗透率提升切换到DDR5内部代际升级（Gen-3/Gen-4），高毛利新品占比提升是净利润超预期核心来源\par}
\paragraph{MS} 澜起科技2Q26初步营收19亿元（环比+28\%，同比+33\%），净利润中位数12亿元（环比+36\%，同比+82\%），高出预期50\%。扣非净利润7亿元（环比+24\%，同比+27\%）。AI推理与智能体任务持续拉动RCD及新型互联芯片需求，PCIe业务在中国AI服务器市场具有独特地位。韩国反垄断调查未影响运营，三季度需求展望强劲。
{\small\color{refgray}数据：2Q26营收19亿元，环比+28\%，同比+33\%\par}
{\small\color{refgray}数据：净利润中位数12亿元，环比+36\%，同比+82\%\par}
{\small\color{refgray}数据：高出MS预期50\%\par}
{\small\color{refgray}数据：扣非净利润7亿元，环比+24\%，同比+27\%\par}
{\small\color{refgray}边际变化：AI需求同时拉动RCD和新互联芯片两条产品线，扣非净利润仍有27\%同比增长，显示可持续性\par}
\paragraph{JPM} 韩国电池行业从抢份额转向算利润，LG新能源获谷歌2.9GWh储能订单验证LFP技术路径，但利润拐点需跟踪封装产能爬坡（下半年出货量从10.6GWh增至19.4GWh）。韩国AI储能招标640MWh，三星SDI获66\%份额。资产出清速度决定哪家公司率先走出估值低谷。
{\small\color{refgray}数据：LG新能源获谷歌2.9GWh储能订单\par}
{\small\color{refgray}数据：韩国AI储能招标640MWh，三星SDI获420MWh（66\%）\par}
{\small\color{refgray}数据：LG新能源下半年出货量预计从10.6GWh增至19.4GWh\par}
{\small\color{refgray}数据：预计2026年四季度实现扣除AMPC后盈亏平衡\par}
{\small\color{refgray}边际变化：市场焦点从订单规模转向利润拐点与闲置资产盘活，封装产能爬坡成为关键变量\par}
\paragraph{UBS} Momenta作为中国最大独立自动驾驶软件商，与24家全球车企合作，覆盖前十大中9家，累计超170个定点项目，68个已量产。独立城市NOA方案市占率64.5\%。预计年安装量从2025年5.1万辆增至2028年250万辆，收入CAGR约50\%，2028年有望盈亏平衡。竞争来自华为、地平线等，车企压价与芯片涨价是风险。
{\small\color{refgray}数据：合作24家全球车企，覆盖前十大中9家\par}
{\small\color{refgray}数据：累计超170个定点项目，68个已量产\par}
{\small\color{refgray}数据：独立城市NOA方案市占率64.5\%\par}
{\small\color{refgray}数据：年安装量从2025年5.1万辆增至2028年250万辆\par}
{\small\color{refgray}边际变化：Momenta从技术供应商演变为主流车企标配，跨品牌数据积累与工程经验成为护城河\par}
\subsubsection*{关键数据}
\begin{itemize}
  \item 3Q26 DRAM ASP环比+21\%（BofA）
  \item ASML 2027年Low-NA EUV产能接近售罄，2028年订单已流入（Citi）
  \item 800VDC预计2030年覆盖79\%新增数据中心容量（Citi）
  \item 中国AI芯片2025年出货约400万颗，供需缺口至少至2027年（DB）
  \item 澜起科技2Q26净利润超预期48\%，DDR5 Gen-3/Gen-4驱动（GS）
  \item 光互联市场2025-2030年CAGR 31.6\%，中际旭创市占率21.2\%（Citi）
  \item Momenta独立城市NOA方案市占率64.5\%，年安装量预计2028年达250万辆（UBS）
  \item LG新能源下半年出货量从10.6GWh增至19.4GWh（JPM）
\end{itemize}
\begin{figure}[htbp]
\centering
\includegraphics[width=0.88\linewidth]{figures/fig_001.jpg}
\caption{Exhibit 2 - Exhibit 3: TrendForce forecasts NAND ASP to normalize to \textbackslash{}\textasciitilde{}+10-15\% QoQ in 3Q26 and \textbackslash{}\textasciitilde{}+0-5\% QoQ in 4Q26 after reaching peak levels of \textbackslash{}\textasciitilde{}85-90\% QoQ and \textbackslash{}\textasciitilde{}55-60\% QoQ in 1Q/2Q26 TrendForce NAND ASP forecast updated as of 30 \$\textasciicircum{}\{th\}\$ Jun}
\end{figure}
\begin{figure}[htbp]
\centering
\includegraphics[width=0.88\linewidth]{figures/fig_007.jpg}
\caption{Figure 1 - \# Nvidia's 800VDC Architecture: Four Generations, from Multi-Stage AC to a Single Integrated DC Conversion NVIDIA's reference architecture defines 4 generations of data center power distribution, each representing a step-change reduction in conversion stages (}
\end{figure}
\begin{figure}[htbp]
\centering
\includegraphics[width=0.88\linewidth]{figures/fig_017.jpg}
\caption{Figure 2 - \#\# Business overview Zhongji Innolight, founded in 2008, is now a world leader in development of optical transceivers. It provides a wide range of high-speed optical solutions for optical communications networking, especially for AI and data center application}
\end{figure}
\begin{figure}[htbp]
\centering
\includegraphics[width=0.88\linewidth]{figures/fig_074.jpg}
\caption{Figure 3 - Figure 5: Momenta leads in key commercial milestones including nominations and SOPs for urban NOA}
\end{figure}
{\small\color{refgray}本节综合 9 篇研究；涉及 BofA、Citi、DB、GS、JPM、MS、UBS。}
\newpage
\subsection{消费与汽车：结构分化与品牌韧性}
\textbf{中国新能源车市场折扣分化，新势力收敛而传统品牌扩大；饮料行业无糖茶取代包装水成为增长引擎；运动服饰折扣纪律与库存压力并存；欧洲卡车需求修复，北美订单增长但需区分补库与真需求；奢侈品品牌复苏逻辑持续验证。}
\subsubsection*{主流共识}
\begin{itemize}
  \item 中国新能源车市场整体折扣率趋于稳定，但品牌间分化加剧，新势力折扣收敛而传统自主品牌仍在扩大让利。
  \item 无糖茶品类已成为饮料行业增长主线，东方树叶等品牌增速远超含糖茶，推动品类结构切换。
  \item 运动服饰行业折扣加深，但头部品牌通过折扣纪律维护品牌力，库存压力成为短期关注点。
\end{itemize}
\subsubsection*{分歧与条件差异}
\begin{itemize}
  \item 新能源车定价权回归信号：新势力折扣收敛是否意味着定价权回归，还是仅为阶段性平衡点，市场存在分歧。
  \item 北美卡车订单增长性质：122\%的订单增长是终端需求复苏还是渠道补库存，决定周期高度。
  \item 奢侈品品牌复苏可持续性：市场对下半年比较基数担忧与品牌连续四个季度加速的证据之间存在分歧。
\end{itemize}
\subsubsection*{机构观点}
\paragraph{Citi} 7月上半月中国新能源乘用车加权平均零售折扣率5.9\%，环比仅微增0.1个百分点，价格调整从全面展开转向局部调整。新势力阵营折扣收敛：理想折扣从4.6\%收窄至4.1\%，特斯拉维持0.6\%，蔚来、小鹏变动有限。传统自主品牌折扣仍在爬升：比亚迪从6.7\%升至6.9\%，吉利折扣最大。豪华品牌价格承压，行业价格竞争未明显升级但品牌间差异显著。
{\small\color{refgray}数据：7月上半月新能源车加权平均零售折扣率5.9\%，环比+0.1ppt\par}
{\small\color{refgray}数据：理想折扣率从4.6\%降至4.1\%，环比-0.6ppt\par}
{\small\color{refgray}数据：比亚迪折扣率从6.7\%升至6.9\%，环比+0.2ppt\par}
{\small\color{refgray}数据：特斯拉折扣率连续两个月维持0.6\%\par}
{\small\color{refgray}边际变化：行业价格调整从全面展开转向局部调整，新势力折扣收敛而传统自主品牌仍在扩大让利。\par}
\paragraph{DB} 小鹏Mona L03 SUV最终定价12.4万至15.7万元，比预售价低约1万元，公布后7分钟内获2万个不可取消订单。设计由前法拉利设计师团队操刀，外观与内饰质感突出。预计2026年月交付约1.5万辆，2027年全年销量达15万辆，延续Mona M03月销过万的成功模式。
{\small\color{refgray}数据：Mona L03定价12.4万-15.7万元，比预售价低约1万元\par}
{\small\color{refgray}数据：公布后7分钟内获2万个不可取消订单\par}
{\small\color{refgray}数据：预计2026年月交付约1.5万辆，2027年全年销量15万辆\par}
{\small\color{refgray}数据：Mona M03月均销量超1万辆\par}
{\small\color{refgray}边际变化：小鹏通过低价SUV定价策略和设计差异化，在主流SUV市场寻求规模增长，订单反应超预期。\par}
\paragraph{GS} 农夫山泉2026年上半年营收预计293亿元，同比增14\%，净利润89亿元，利润率从29.7\%升至30.4\%。茶饮料取代包装水成为最大增长引擎，收入占比从2022年的21\%升至43\%，包装水从55\%降至33\%。东方树叶同比增35\%，茶π仅增3\%，无糖茶线下增速16.1\%领先其他品类。包装水同比增6\%，但尚未回到1H23水平。
{\small\color{refgray}数据：1H26营收预计293亿元，同比+14\%\par}
{\small\color{refgray}数据：净利润89亿元，利润率30.4\%\par}
{\small\color{refgray}数据：茶饮料收入占比从2022年21\%升至43\%\par}
{\small\color{refgray}数据：东方树叶同比+35\%，茶π仅+3\%\par}
{\small\color{refgray}边际变化：茶饮料取代包装水成为最大增长引擎和利润来源，无糖茶品类增速远超含糖茶，品类结构切换加速。\par}
\paragraph{GS} 特步核心品牌二季度零售额同比下滑中个位数，索康尼增速从20\%+收窄至低个位数。行业折扣加深但特步折扣维持在25-30\%同比持平，线上折扣改善。渠道库存从4.0-4.5个月升至4.5-5.0个月，库存上升而折扣未加深。下调2026-2028年收入预测2-3\%，净利润调低1-3\%，维持毛利率受折扣纪律支撑判断。
{\small\color{refgray}数据：核心品牌Q2零售额同比下滑中个位数\par}
{\small\color{refgray}数据：索康尼增速从Q1的20\%+降至低个位数\par}
{\small\color{refgray}数据：折扣维持25-30\%，同比持平\par}
{\small\color{refgray}数据：渠道库存从4.0-4.5个月升至4.5-5.0个月\par}
{\small\color{refgray}边际变化：行业折扣加深环境下特步选择克制策略，库存上升但折扣未加深，短期周转效率换取品牌定价权。\par}
\paragraph{MS} Burberry第一财季同店销售增长5\%，在更难比较基数上实现500bps同比加速，但报价当天跑输板块约500bps。MS认为市场反应过度：品牌已连续四个季度同店销售同比加速，增长质量改善——手袋品类转正、欧洲集群转正、全价销售占比提升、折扣同比减少。H1毛利率微降是库存拨备正常化而非折扣问题，全年毛利率仍扩张。
{\small\color{refgray}数据：Q1同店销售同比+5\%，较去年同期加速500bps\par}
{\small\color{refgray}数据：连续四个季度同店销售同比加速\par}
{\small\color{refgray}数据：手袋品类转正，欧洲集群转正\par}
{\small\color{refgray}数据：H1毛利率微降因库存拨备正常化，全年仍扩张\par}
{\small\color{refgray}边际变化：市场对下半年比较基数担忧过度，品牌复苏逻辑持续验证，增长质量改善而非一次性反弹。\par}
\paragraph{MS} 6月香槟出口量同比增8\%，二季度出口增速达6.3\%，扭转4月持平、5月下滑1\%走势。LVMH约60\%香槟销量来自出口，在出口市场量份额约34\%、价值份额约50\%。凯歌品牌在美国吸收其全球产量约40\%。当前市场预期二季度香槟与葡萄酒板块同比+1.4\%，出口加速带来小幅超预期可能，但该业务占集团营收不足4\%。
{\small\color{refgray}数据：6月香槟出口量同比+8\%\par}
{\small\color{refgray}数据：二季度出口增速6.3\%\par}
{\small\color{refgray}数据：LVMH约60\%香槟销量来自出口\par}
{\small\color{refgray}数据：LVMH出口市场量份额34\%，价值份额50\%\par}
{\small\color{refgray}边际变化：出口加速为LVMH香槟业务带来小幅超预期可能，但整体行业仍处于调整阶段。\par}
\paragraph{MS} 沃尔沃二季度调整后营业利润148亿瑞典克朗，同比增10\%，利润率11.7\%。卡车业务利润超预期8\%，订单增长33\%，北美订单增122\%。管理层上调欧洲卡车市场预期至31.5万辆，中国上调至88万辆，北美维持26.5万辆。北美订单增长更多来自补库存而非终端需求，货运市场仍处供给驱动阶段。
{\small\color{refgray}数据：调整后营业利润148亿瑞典克朗，同比+10\%\par}
{\small\color{refgray}数据：利润率11.7\%，同比+70bps\par}
{\small\color{refgray}数据：集团订单增长33\%，北美订单+122\%\par}
{\small\color{refgray}数据：欧洲卡车市场预期上调至31.5万辆\par}
{\small\color{refgray}边际变化：欧洲和中国需求修复验证，但北美订单增长主要来自补库存，货运市场尚未进入需求驱动阶段。\par}
\subsubsection*{关键数据}
\begin{itemize}
  \item 7月上半月新能源车加权平均零售折扣率5.9\%，环比+0.1ppt
  \item 理想折扣率从4.6\%降至4.1\%，环比-0.6ppt
  \item 小鹏Mona L03定价12.4万-15.7万元，7分钟获2万订单
  \item 农夫山泉1H26茶饮料收入占比升至43\%，东方树叶+35\%
  \item 特步核心品牌Q2零售额同比下滑中个位数，库存升至4.5-5.0个月
  \item Burberry连续四个季度同店销售同比加速，Q1+5\%
  \item 6月香槟出口量同比+8\%，LVMH出口市场价值份额50\%
  \item 沃尔沃Q2订单增长33\%，北美+122\%，欧洲市场预期上调至31.5万辆
\end{itemize}
\begin{figure}[htbp]
\centering
\includegraphics[width=0.88\linewidth]{figures/fig_012.jpg}
\caption{Figure 1 - Figure 1. NEV-PV retail discounts by OEMs © 2026 Citi Inc. No redistribution without Citi's written permission. Figure 2. Sector weighted average NEV-PV discount vs MSRP © 2026 Citi Inc. No redistribution without Citi's written permission. Figure 3. NEV-PV ret}
\end{figure}
\begin{figure}[htbp]
\centering
\includegraphics[width=0.88\linewidth]{figures/fig_046.jpg}
\caption{Exhibit 1 - Exhibit 1: Champagne Shipments Growth YoY (in volume) Exhibit 2: Champagne Exports Growth YoY (in volume) \#\# Additional Charts}
\end{figure}
\begin{figure}[htbp]
\centering
\includegraphics[width=0.88\linewidth]{figures/fig_056.jpg}
\caption{Exhibit 4 - Exhibit 4: Group Adjusted EBIT Margin (\%)}
\end{figure}
\begin{figure}[htbp]
\centering
\includegraphics[width=0.88\linewidth]{figures/fig_013.jpg}
\caption{Figure 4 - © 2026 Citi Inc. No redistribution without Citi's written permission. © 2026 Citi Inc. No redistribution without Citi's written permission. Figure 6. CPCA: Monthly Number of Car Model Adopted Price Cut © 2026 Citi Inc. No redistribution without Citi's written }
\end{figure}
{\small\color{refgray}本节综合 7 篇研究；涉及 Citi、DB、GS、MS。}
\newpage
\subsection{医疗健康：结构性分化与预期校准}
\textbf{医疗健康板块正经历结构性分化：美国管理式医疗龙头超预期但股价反应平淡，反映市场对盈利可持续性的定价已前置；手术机器人护城河受侵蚀，中国市场断崖式下滑与耗材业务承压；中国生物科技板块估值修复主要由资金轮动驱动，持续性待观察；新药放量速度成为药企分水岭，部分新药需数倍于历史标杆的放量速度才能达成共识预期。}
\subsubsection*{主流共识}
\begin{itemize}
  \item 美国管理式医疗龙头联合健康二季度业绩全面超预期，营收1120亿美元，调整后EPS 6.38美元，均高于市场预期。
  \item 中国生物科技板块估值在6月底触及低点后快速反弹，P/S从3.0倍回升至3.8倍，流动性占比从6\%升至9.7\%。
  \item 新药上市后的处方量放量速度是判断药企营收预期能否兑现的核心指标，MS为四款新药建立了可追踪的处方量门槛框架。
  \item 手术机器人龙头直觉外科美国手术量增速放缓至12\%，且维持全年指引未上调，打破其长期beat and raise记录。
\end{itemize}
\subsubsection*{分歧与条件差异}
\begin{itemize}
  \item 联合健康商业险利润率恢复时间线拉长至2027年之后，DB认为这是“延迟而非挫折”，未来可通过定价调整和AI驱动效率提升修复；但市场对此反应平淡，股价高开低走。
  \item GSK慢性咳嗽药物camlipixant两项III期试验一胜一负，资产终止开发，中期营收缺口固化；DB此前曾乐观预期该资产可能改变叙事，实际结果低于预期。
  \item 直觉外科中国市场二季度仅安装2台系统（去年同期29台），本土竞争加剧；但印度市场dV5获批，高端产品在价格敏感市场推广存在挑战。
  \item 中国生物科技板块估值修复由资金轮动和政策预期驱动，而非基本面改善；MS指出其持续性取决于后续政策落地和交易进展。
\end{itemize}
\subsubsection*{机构观点}
\paragraph{DB} 联合健康二季度营收1120亿美元、调整后EPS 6.38美元均超预期，全年指引上调至19.50-20.00美元，MLR指引下调约70bp至88.1\%。但股价高开低走，市场认为利好已定价。商业险利润率恢复推迟至2027年之后，DB定性为“延迟而非挫折”，未来可通过定价调整和AI驱动的效率提升修复。OptumHealth利润率从2.6\%升至5.0\%，恢复超预期。
{\small\color{refgray}数据：营收1120亿美元，高于预期1106亿\par}
{\small\color{refgray}数据：调整后EPS 6.38美元，高于预期4.89美元\par}
{\small\color{refgray}数据：全年EPS指引上调至19.50-20.00美元\par}
{\small\color{refgray}数据：MLR指引下调约70bp至88.1\%\par}
{\small\color{refgray}边际变化：市场对超预期业绩反应平淡，表明盈利可持续性已被定价；商业险利润率恢复时间线拉长，盈利改善节奏更多依赖Optum。\par}
\paragraph{DB} GSK慢性咳嗽药物camlipixant两项III期试验CALM-1达标、CALM-2未达标，资产终止开发。DB此前估值每股74便士（基于2037年峰值销售8亿英镑）需归零。GSK 2031年营收目标超400亿英镑，DB和共识预估分别为350亿和370亿英镑，camlipixant失败使缺口固化。慢性咳嗽领域至今无FDA批准的P2X3药物。
{\small\color{refgray}数据：CALM-1入组825人，CALM-2入组975人\par}
{\small\color{refgray}数据：camlipixant NPV每股74便士，峰值销售假设8亿英镑（2037E）\par}
{\small\color{refgray}数据：GSK 2031年营收目标超400亿英镑，DB预估350亿英镑\par}
{\small\color{refgray}数据：共识预期2031年camlipixant销售仅4.15亿英镑\par}
{\small\color{refgray}边际变化：DB此前预期camlipixant可能带来正面叙事，实际结果一胜一负且资产终止，中期营收缺口从“可能被填补”变为“大概率维持”。\par}
\paragraph{DB} 直觉外科二季度全球手术量增长15\%符合预期，但美国仅增12\%低于预期，且维持全年指引未上调。中国市场二季度仅安装2台系统（上半年6台，去年同期29台），本土竞争挤压中标率。Extended Use 2.0预计2027年上半年推出，可能拉低单次手术耗材收入约7\%。三重护城河（美国手术量增速、中国增长引擎、核心耗材业务）正同时变薄。
{\small\color{refgray}数据：美国手术量增速12\%，低于市场预期\par}
{\small\color{refgray}数据：中国二季度仅安装2台系统，上半年6台（去年同期29台）\par}
{\small\color{refgray}数据：全年美国手术量增速预测下调110bp至12.1\%\par}
{\small\color{refgray}数据：Extended Use 2.0预计拉低单次手术耗材收入约7\%\par}
{\small\color{refgray}边际变化：公司历史上首次在手术量上“不调高”，维持指引已是隐性下调；中国业务从增长引擎变为拖累，本土竞争加剧。\par}
\paragraph{MS} 美国药品市场总处方量同比下降5.3\%，过去12周滚动增速仅1.2\%。四款新药放量速度分化：BMY的Cobenfy周处方约3300张，要达到2026年2.89亿美元共识收入需2-5倍于同类药物放量速度；GILD的Yeztugo周处方约1750张，注射剂型增长明显；VRTX的Journavx周处方约23460张，医院处方未被统计；JNJ的Icotyde上市3周周处方约800张，医生反馈积极。
{\small\color{refgray}数据：全美总处方量同比下降5.3\%\par}
{\small\color{refgray}数据：Cobenfy周处方约3300张，需2-5倍于同类放量速度\par}
{\small\color{refgray}数据：Yeztugo周处方约1750张，注射剂型从750增至1030张\par}
{\small\color{refgray}数据：Journavx周处方约23460张，平均用药时长10-11天\par}
{\small\color{refgray}边际变化：新药放量速度成为药企分水岭，MS建立可复用的处方量门槛框架，将共识预期拆解为可追踪的周处方量轨迹。\par}
\paragraph{MS} 中国生物科技板块P/S从5月底3.5倍压缩至6月下旬3.0倍，随后反弹至7月中旬3.8倍；EV/S从2.8倍降至2.4倍后回升至3.2倍。流动性占比从6\%升至9.7\%。估值修复主要由期末资金再平衡和科技板块资金轮动驱动，叠加基药目录政策利好和BINSA悬疑解除。MS指出，修复持续性取决于后续政策落地和交易进展。
{\small\color{refgray}数据：P/S从3.0倍反弹至3.8倍，反弹幅度约27\%\par}
{\small\color{refgray}数据：EV/S从2.4倍回升至3.2倍\par}
{\small\color{refgray}数据：流动性占比从6\%升至9.7\%\par}
{\small\color{refgray}数据：ASCO期间流动性一度升至7.5\%\par}
{\small\color{refgray}边际变化：估值修复由资金轮动和政策预期驱动，而非基本面改善；快速反弹（两周内27\%）伴随交易性资金涌入，持续性待观察。\par}
\subsubsection*{关键数据}
\begin{itemize}
  \item 联合健康营收1120亿美元，调整后EPS 6.38美元，均超预期
  \item 联合健康MLR指引下调约70bp至88.1\%
  \item GSK camlipixant NPV每股74便士归零，中期营收缺口固化
  \item 直觉外科中国二季度仅安装2台系统（去年同期29台）
  \item 直觉外科美国手术量增速12\%，低于预期
  \item Cobenfy周处方约3300张，需2-5倍于同类放量速度
  \item 中国生物科技P/S从3.0倍反弹至3.8倍，流动性占比升至9.7\%
  \item 全美总处方量同比下降5.3\%
\end{itemize}
\begin{figure}[htbp]
\centering
\includegraphics[width=0.88\linewidth]{figures/fig_039.jpg}
\caption{Exhibit 37 - Exhibit 1: MS-Estimated Cobenfy TRx trajectory to achieve 2025/2026 consensus estimates vs. schizophrenia launch analogs Exhibit 2: Cobenfy Launch vs. Recent antipsychotic launches (all indications) Exhibit 3: MS-Estimated Jour}
\end{figure}
\begin{figure}[htbp]
\centering
\includegraphics[width=0.88\linewidth]{figures/fig_044.jpg}
\caption{Exhibit 1 - Exhibit 1: China Biotech 2030e Price-to-Sales Ratio mid-Jul 2026 Exhibit 2: China Biotech 18A\textbackslash{}* Liquidity – mid-Jul 2026 MS ASIA LIMITED+ Jack Lin}
\end{figure}
\begin{figure}[htbp]
\centering
\includegraphics[width=0.88\linewidth]{figures/fig_022.jpg}
\caption{Figure 1 - Figure 1: Camlipixant sales DBe vs consensus (GBPm)}
\end{figure}
\begin{figure}[htbp]
\centering
\includegraphics[width=0.88\linewidth]{figures/fig_023.jpg}
\caption{Figure 4 - Figure 4: Chronic cough clinical data}
\end{figure}
{\small\color{refgray}本节综合 5 篇研究；涉及 DB、MS。}
\newpage
\subsection{区域市场与宏观：增长分化与政策预期调整}
\textbf{7月全球宏观呈现结构性分化：中国增长动能从全面走弱转为出口扰动与财政加速的组合，地产量缩正向价跌传导；欧洲电气化政策信号拉长公用事业投资确定性窗口；韩国杠杆ETF市场在监管加码前已自发去杠杆。锂电产业链则面临库存跑赢销量的主动去库压力。}
\subsubsection*{主流共识}
\begin{itemize}
  \item 中国7月出口量增速因台风扰动阶段性放缓，但运价上涨反映供给端紧张未缓解。
  \item 欧洲电气化政策目标（2040年电力需求约5000 TWh）虽激进，但基于成熟技术，可行性高于纯脱碳目标。
  \item 韩国单股杠杆ETF新规（最低保证金提至3000万韩元）将显著抑制散户参与，但市场已先于政策调整。
\end{itemize}
\subsubsection*{分歧与条件差异}
\begin{itemize}
  \item 中国地产市场：MS认为二手房销售减速将加速房价调整，尤其二线及以下城市；但一线城市新房库存去化与二手房价格环比仍正增长，显示城市层级分化。
  \item 锂电产业链：JEF调研显示库存增速连续超过销量增速，指向主动去库与价格承压；但MS周报显示7月锂电池排产创新高且储能占比突破40\%，需求结构改善可能缓解过剩压力。
  \item 中国财政节奏：JPM高频数据提示7月政府债发行已超6月全月，但MS月度追踪显示新房销售与土地市场仍弱，财政发力能否有效传导至实体经济存在分歧。
\end{itemize}
\subsubsection*{机构观点}
\paragraph{JPM} 7月前两周中国宏观信号分化：出口量增速因台风扰动从6月17.2\%降至8.7\%，但集装箱运价继续上涨（美西+14.9\%），供给紧张未缓解；政府债券发行1.17万亿元已超6月全月，财政加速落地；工业生产中炼油改善（沥青开工率回升）但汽车与钢铁承压。二季度GDP不及预期后，市场对下半年增长节奏重新评估，货币宽松预期从Q3推迟至Q4。
{\small\color{refgray}数据：7月前两周离港船舶载重吨同比+8.7\%（6月+17.2\%）\par}
{\small\color{refgray}数据：CCFI美西航线较两周前上涨14.9\%\par}
{\small\color{refgray}数据：7月政府债发行1.17万亿元，已超6月全月\par}
{\small\color{refgray}数据：7天逆回购净投放1970亿元但回笼300亿元\par}
{\small\color{refgray}边际变化：出口增速放缓主因台风扰动而非需求走弱，财政发债提速成为新支撑；货币宽松预期从Q3推迟至Q4。\par}
\paragraph{JEF} 三季度锂电产业链动能明显走弱：受访企业预期锂盐价格将环比下降7-8\%，新订单与积压订单扩散指数跌破50\%荣枯线，电池生产商订单减速尤为显著。库存同比增速已连续超过销量增速，过去三年仅出现两次，通常对应价格承压阶段。超过50\%受访企业认为下游客户库存偏高，后续补库意愿减弱。
{\small\color{refgray}数据：锂盐价格预期环比下降7-8\%\par}
{\small\color{refgray}数据：积压订单扩散指数跌破50\%\par}
{\small\color{refgray}数据：库存同比增速连续超过销量增速（过去三年仅两次）\par}
{\small\color{refgray}数据：超50\%受访企业认为下游客户库存偏高\par}
{\small\color{refgray}边际变化：从二季度持平预期转为明确的下行预期，主动去库阶段确认。\par}
\paragraph{MS} 欧盟电气化计划草案若落地，2040年电力需求将达约5000 TWh，较当前预测高约50\%，即使仅实现一半也足以支撑公用事业进入新增长时代。欧盟电气化率仅23\%且十年未提升，但追赶空间明确。电网与清洁发电资本开支周期拉长至2040年，资产估值需重新审视。
{\small\color{refgray}数据：欧盟2040年电力需求目标约5000 TWh（MSe为3400 TWh）\par}
{\small\color{refgray}数据：当前电气化率仅23\%，中国/日本接近30\%\par}
{\small\color{refgray}数据：发电装机需从1132 GW增至约3000 GW\par}
{\small\color{refgray}数据：计划到2040年减少2600亿欧元化石燃料进口成本\par}
{\small\color{refgray}边际变化：政策信号将投资确定性窗口从2030年延伸至2040年，公用事业估值逻辑从短期周期转向长期增长。\par}
\paragraph{MS} 6月中国新房销售面积同比-7\%（5月+1\%），二手房增速从+14\%收窄至+6\%，实时高频显示进一步放缓至+8\%。挂牌量创新高城市占比从31\%升至35\%，供给端不确定性未见顶。新房库存月数升至31.8个月，三四线达42.7个月。土地成交面积同比-23\%，开发商新推货量同比-14\%，下半年供应节奏将受影响。
{\small\color{refgray}数据：CREIS 65城新房销售面积同比-7\%（5月+1\%）\par}
{\small\color{refgray}数据：33城二手房销售增速从+14\%降至+6\%\par}
{\small\color{refgray}数据：70城新房库存月数31.8个月（5月31.6个月）\par}
{\small\color{refgray}数据：挂牌量创新高城市占比35\%（5月31\%）\par}
{\small\color{refgray}边际变化：从量缩加速传导至价跌，二手房销售减速将加速房价调整，尤其二线及以下城市。\par}
\paragraph{UBS} 韩国单股杠杆ETF新规（最低保证金从300万提至3000万韩元、暂停新产品上市）旨在抑制散户投机，但市场已先于政策调整：AUM从6月25日峰值24万亿韩元降至17万亿韩元（-29\%）。杠杆ETF亏损远超底层资产（SK海力士自5月27日跌7\%，对应2倍ETF跌32\%），负复利效应显著。新规对个人投资者效果最明显，相当于第三收入阶层家庭金融资产的27\%。
{\small\color{refgray}数据：单股杠杆ETF AUM从24万亿韩元降至17万亿韩元（-29\%）\par}
{\small\color{refgray}数据：最低保证金从300万韩元提至3000万韩元\par}
{\small\color{refgray}数据：SK海力士自5月27日跌7\%，对应2倍ETF跌32\%\par}
{\small\color{refgray}数据：新规影响相当于第三收入阶层家庭金融资产的27\%\par}
{\small\color{refgray}边际变化：市场自发去杠杆已先于监管完成大部分调整，新规进一步抑制散户参与但边际影响递减。\par}
\paragraph{MS} 7月锂电池排产创历史新高，储能电池需求占比首次突破40\%，改变了过去依赖电动车的季节性特征。但锂盐价格继续调整（电池级碳酸锂周跌4.7\%，氢氧化锂周跌5.1\%），产量新高与价格新低并存，反映产能释放快于需求消化。铜铝库存下降（沪铜库存-34.9\%，沪铝库存-1.1\%），对价格形成支撑。
{\small\color{refgray}数据：7月锂电池排产创历史新高\par}
{\small\color{refgray}数据：储能电池需求占比首次突破40\%\par}
{\small\color{refgray}数据：电池级碳酸锂周跌4.7\%，氢氧化锂周跌5.1\%\par}
{\small\color{refgray}数据：沪铜库存周降34.9\%，沪铝库存周降1.1\%\par}
{\small\color{refgray}边际变化：储能占比跃升使锂电池需求从单一依赖电动车转向双轮驱动，但价格竞争与放量并存。\par}
\subsubsection*{关键数据}
\begin{itemize}
  \item 中国7月前两周出口量增速从6月17.2\%降至8.7\%（台风扰动）
  \item 韩国单股杠杆ETF AUM从峰值24万亿韩元降至17万亿韩元（-29\%）
  \item 欧盟2040年电力需求目标约5000 TWh，较MSe高约50\%
  \item 中国6月新房销售面积同比-7\%（5月+1\%），二手房增速从+14\%降至+6\%
  \item 7月锂电池排产创新高，储能占比首次突破40\%
  \item 锂电产业链库存增速连续超过销量增速（过去三年仅两次）
  \item 70城新房库存月数31.8个月，三四线42.7个月
  \item 沪铜库存周降34.9\%
\end{itemize}
\begin{figure}[htbp]
\centering
\includegraphics[width=0.88\linewidth]{figures/fig_035.jpg}
\caption{Figure 1 - \textbackslash{}- Port tracking suggests that export volume growth may have moderated mtd in July, as departing shipments deadweight tonnage (excluding tankers) rose 8.7\% oya mtd in July (vs. 17.2\% in June). This reflects shipping disruption caused by Typhoon Bavi, and port }
\end{figure}
\begin{figure}[htbp]
\centering
\includegraphics[width=0.88\linewidth]{figures/fig_053.jpg}
\caption{Exhibit 2 - Exhibit 5: The EU's 46\% target would imply EU power demand of 5,206 TWh, a 4.9\% CAGR from 2026, vs MSe 3,400 TWh and 2.1\% CAGR}
\end{figure}
\begin{figure}[htbp]
\centering
\includegraphics[width=0.88\linewidth]{figures/fig_062.jpg}
\caption{Exhibit 1 - Exhibit 1: Monthly unit sales – Tier 1 cities Exhibit 2: Monthly unit sales – Tier 2 cities Exhibit 3: Monthly unit sales – Tier 3 cities Exhibit 4: Monthly unit sales – Total 65 cities}
\end{figure}
\begin{figure}[htbp]
\centering
\includegraphics[width=0.88\linewidth]{figures/fig_066.jpg}
\caption{Figure 3 - Figure 5: AUM of single-stock (SEC, SKH) leveraged ETFs listed overseas}
\end{figure}
{\small\color{refgray}本节综合 6 篇研究；涉及 JPM、JEF、MS、UBS。}
\newpage
\section{辅助信号｜战略咨询}
本板块整合波士顿咨询三份报告，聚焦AI在政府信任、国防维护与军工制造中的战略应用，揭示技术落地的信任前提、效率瓶颈与流程重构路径。
\subsection{AI落地的信任前提与场景选择}
\textbf{公民对政府使用生成式AI的接受度高于预期，但前提是政府先解决信任底层问题；低风险、高透明度的场景是切入关键。}
\begin{itemize}
  \item 全球数字政府服务净满意度65\%，两年仅微增2个百分点，但75\%的公民期待政府服务达到顶级私企水平，期望增速远超改进速度。
  \item 63\%的受访者愿意通过GenAI获取简单政府服务，71\%接受多语言沟通，69\%支持政府雇员使用AI辅助工具；抽象焦虑与具体接受之间存在清晰分界线。
  \item 在“用GenAI自动化服务准入决策”和“监控公众情绪”等高不确定性场景上，各司法管辖区接受度差异显著，政府需根据本地社会许可边界设计用例。
  \item 报告提出三步走策略：先小规模高透明度实践积累信任，再用制度和规则锁定信任，最后才有资格讨论大规模创新——顺序不能颠倒。
\end{itemize}
\begin{figure}[htbp]
\centering
\includegraphics[width=0.88\linewidth]{figures/fig_112.jpg}
\caption{Exhibit 1 - Exhibit 1: We surveyed 41,600 regular internet users across 48 jurisdictions \$\textasciicircum{}\{1\}\$ Special Administration Region of the People's Republic of China In conducting this research, we also held interviews and gathered qualitative inp}
\end{figure}
\subsection{AI重构复杂装备维护与供应链效率}
\textbf{全球空军战机任务出勤率长期低于70\%，问题不在投入规模而在流程效率；AI通过知识整合、供应链可视化和流程质疑，可显著提升战备率。}
\begin{itemize}
  \item 美、法、德战机任务可用率常规性低于70\%，英国2024年F-35机队仅三分之一具备完全任务能力；合格维护人员不足是直接原因。
  \item GenAI可将分散的纸质手册、技术文档整合为交互式界面，使经验不足的维护人员快速定位故障、列出所需工具和备件，将“十年经验”压缩为一次查询。
  \item 备件短缺并非真实痛点，而是备件“在某个角落但找不到”；AI供应链控制塔可实时追踪库存位置，将停飞原因从“缺件”转为“找件”。
  \item 报告提出“可靠性控制委员会”机制：用AI识别导致大量停飞天数的长尾零件和流程，反向推动供应商增产、设计变更，甚至质疑既有检查流程的合理性。
\end{itemize}
\subsection{欧洲军工产能瓶颈：流程效率优先于资本开支}
\textbf{欧洲军工需求激增但制造端规模化能力不足，核心瓶颈在于流程低效而非厂房面积；先做流程诊断再决定扩产，是避免资本浪费的关键。}
\begin{itemize}
  \item 欧洲军工企业多为中小企业，不愿自掏腰包研究产能，担心需求节奏变化导致资产闲置；真正瓶颈不在需求端，而在制造端的规模化能力。
  \item 报告指出，现有设施产能受限常源于流程中的低效环节，如设计到生产的转换过程、测试阶段的时间消耗，这些并不需要新增资本开支就能解决。
  \item 供应链脆弱性来自长距离依赖和地缘优先权：海外供应商在地缘紧张时会优先服务本国OEM，欧洲客户只能排在后面；应对思路是“选择性本地化”。
  \item 隐含逻辑：在投入真金白银之前，先问现有产线是否已跑到了极限——顺序错了，钱可能白花。
\end{itemize}
\newpage
\section{辅助信号｜智库/国际机构}
本板块整合经合组织（OECD）与世贸组织（WTO）当日发布的多份研究报告，聚焦三大主题：全球价值链升级路径与制度瓶颈、劳动力市场与生产率的结构性障碍、以及贸易救济与气候融资的新工具。各报告均被实质整合至对应主题中。
\subsection{全球价值链升级：从参与度到附加值的结构性跃迁}
\textbf{经合组织以哥斯达黎加电子业为案例，揭示了嵌入全球价值链的国家向高附加值环节移动的可行路径与量化收益，同时指出制度碎片化是制约升级的关键瓶颈。}
\begin{itemize}
  \item 哥斯达黎加电子业全球价值链参与度较低，且上游度指数自2015年以来持续下降，表明其生产活动正进一步向低附加值的组装环节集中（R030）。
  \item 结构引力模型估算显示，若哥斯达黎加能在电子价值链中向上游移动（如设计、晶圆制造），其出口、产出和实际工资均将出现显著增长（R030）。
  \item 乌兹别克斯坦水安全投资环境评估显示，宏观框架得分较高，但制度碎片化、数据缺口和成本回收能力弱是最大短板，导致PPP项目落地困难（R035）。
  \item 乌兹别克斯坦2019年通过PPP法并设立发展署，但实际投资仍依赖政府预算，根源在于价格信号不真实、风险分配不清晰、数据基础薄弱（R035）。
  \item 两篇报告共同指向：价值链升级不仅需要技术方案，更需要制度环境——包括清晰的职责划分、可验证的数据基础和有效的价格机制。
\end{itemize}
\begin{figure}[htbp]
\centering
\includegraphics[width=0.88\linewidth]{figures/fig_086.jpg}
\caption{Figure 1 - \#\# 4. Empirical application: the case of Costa Rica Costa Rica's trade performance is associated with a strong capacity to attract large inflows of Foreign Direct Investment (FDI), which have been concentrated in sectors such as life sciences, light and advanc}
\end{figure}
\begin{figure}[htbp]
\centering
\includegraphics[width=0.88\linewidth]{figures/fig_107.jpg}
\caption{Figure 3 - \#\# 3.1. An overview of the OECD Scorecard for assessing the enabling environment for investment in water security The purpose of the OECD Scorecard is to pinpoint issues that hamper water security-related investments by public and private investors. The Scorec}
\end{figure}
\subsection{劳动力市场与生产率：竞业限制与AI采纳的结构性分化}
\textbf{经合组织首次用跨国企业数据证实竞业限制条款抑制生产率提升型劳动力再配置和知识扩散；同时，高等教育领域GenAI渗透率虽高，但深度整合远未启动，学生与教师之间存在使用鸿沟。}
\begin{itemize}
  \item 竞业限制条款覆盖率每提高10个百分点，生产率提升型劳动力再配置强度下降约17\%，整体劳动生产率水平下降约1.9\%（R032）。
  \item 竞业限制阻碍了劳动力从低效企业向高效企业的流动，也减缓了落后企业追赶前沿的速度，对宏观生产率构成系统性拖累（R032）。
  \item 英国本科生GenAI使用率从2024年初的66\%飙升至2026年的95\%，但全球仅17\%的学术人员自评为高级或专家级AI用户，88\%的教师使用AI仅停留在最低或中等程度（R033）。
  \item 学生主要用GenAI完成摘要、翻译等常规任务，而非深度学习或解决复杂问题，当前使用模式是对现有流程的效率优化，而非教学范式变革（R033）。
  \item 两篇报告共同揭示：无论是劳动力市场制度还是技术采纳，广泛覆盖不等于深度整合，结构性分化可能削弱长期生产率增长潜力。
\end{itemize}
\begin{figure}[htbp]
\centering
\includegraphics[width=0.88\linewidth]{figures/fig_096.jpg}
\caption{Figure 1 - \#\# Box 1. Non-compete clauses and business entry Non-compete clauses directly prohibit employees from starting competing businesses. Though the practice is often justified by the need for incumbent firms to protect legitimate business interests and trade secre}
\end{figure}
\begin{figure}[htbp]
\centering
\includegraphics[width=0.88\linewidth]{figures/fig_101.jpg}
\caption{Figure 1 - \#\# Nearly all students and most academic staff have experience with GenAI In the three years since the public release of the first of a new generation of generative artificial intelligence (“GenAI”) tools in late 2022, they have developed from a novelty to bei}
\end{figure}
\subsection{贸易救济与气候融资：新工具、新风险与新机遇}
\textbf{美国启动对进口羊肉的保障调查，时间表跨越政治周期，凸显贸易救济工具的政治敏感性；经合组织则提出规模化碳信用模式，以政府为主体推动结构性减排，但环境完整性仍是关键变量。}
\begin{itemize}
  \item 美国国际贸易委员会于2026年7月13日启动对进口羔羊肉的保障措施调查，认定案件异常复杂，最终裁决可能跨越中期选举后的权力交接期（R036）。
  \item 调查不要求证明不公平贸易行为，只需证明进口激增与国内产业严重损害之间的因果关系，门槛最低但政治影响最大（R036）。
  \item 规模化碳信用模式以政府为主体，通过政策执行获得碳信用，截至2025年已有超过80个司法管辖区参与，但仅签发1.54亿个单位，占全球碳信用供给约6\%（R034）。
  \item 规模化碳信用的环境完整性面临基线虚高、额外性存疑、碳泄漏等挑战，其能否兑现承诺取决于能否建立严格的核算与调整机制（R034）。
  \item 两篇报告共同指向：无论是贸易救济还是气候融资，新工具的出现都伴随着制度设计与执行能力的考验，政策制定者需在效率与完整性之间寻求平衡。
\end{itemize}
\begin{figure}[htbp]
\centering
\includegraphics[width=0.88\linewidth]{figures/fig_102.jpg}
\caption{Figure 2 - \textbackslash{}- Article 4 outlines how governments set their mitigation strategies, targets and plans. Scaled-up crediting should support both the ambition and achievement of these efforts, and remain complementary to mitigation that would occur in the absence the creditin}
\end{figure}
\begin{figure}[htbp]
\centering
\includegraphics[width=0.88\linewidth]{figures/fig_103.jpg}
\caption{Figure 2 - \#\# Box 2.2. Understanding corresponding adjustments Corresponding adjustments are an accounting measure under Article 6 to prevent GHG emission reductions or removals from being counted twice. Under Article 6, corresponding adjustments effectively transfer mit}
\end{figure}
\subsection{咨询与企业战略}
\textbf{用于补充企业战略、管理层议题和产业落地视角。}
\begin{itemize}
  \item 全球服装和鞋类行业最广泛使用的工厂认证之一——WRAP（Worldwide Responsible Accredited Production）认证，在经合组织最新发布的评估中，未能通过其尽职调查标准的对齐检验。这份发布于2026年的报告，不是对个别工厂的合规抽查，而是对WRAP...
  \item 当政策挑战日益交叉、资源持续收紧、公众对公共服务的要求不断攀升，政府中枢——即总理府、内阁办公室或类似机构——的角色正在从传统的协调者，转向战略顾问和交付推动者。经合组织最新发布的《建设政府中枢能力以引导和交付复杂优先事项》综合报告，基于对保加利亚、爱沙尼亚、希腊、爱尔兰、波兰和...
\end{itemize}
\begin{figure}[htbp]
\centering
\includegraphics[width=0.88\linewidth]{figures/fig_081.jpg}
\caption{Figure 1 - This section presents the findings of the assessment. Findings are presented in a bar chart displaying the number of criteria found to be fully aligned, partially aligned, or not aligned per step and calculation formulae. A number of strengths and areas for im}
\end{figure}
\begin{figure}[htbp]
\centering
\includegraphics[width=0.88\linewidth]{figures/fig_091.jpg}
\caption{Figure 1 - The remainder of this chapter examines these dimensions in turn, drawing on the findings from the OECD countries that contributed to the project mentioned above. It draws on data collected through surveys undertaken with the six country representatives and pee}
\end{figure}
\subsection{AI、技术与生产率}
\textbf{用于判断 AI 投资周期、生产率扩散和产业约束是否正在改变资产定价。}
\begin{itemize}
  \item 2026年6月，经合组织发布了对公平劳工协会（FLA）制造业认证标准的评估报告。结论是“部分对齐”——一个看似温和、实则指向深层问题的评级。这份报告不是对FLA的否定，而是揭示了当前行业劳工认证体系中最棘手的矛盾：标准要求企业做什么，与标准真正能验证什么，之间存在系统性落差。
\end{itemize}
\begin{figure}[htbp]
\centering
\includegraphics[width=0.88\linewidth]{figures/fig_076.jpg}
\caption{Figure 1 - This chapter presents the findings of the assessment. Findings are presented in a bar chart displaying the number of criteria found to be fully aligned, partially aligned or not aligned per step and calculation formulae. A number of strengths and areas for imp}
\end{figure}
\section{结语}
后续需验证的关键节点包括：存储价格走势是否如TrendForce预测回落、北美卡车订单性质在8月数据中能否明确、中国财政发力向实体经济的传导效果、以及锂电产业链库存去化进度与储能需求持续性。
\newpage
\section{覆盖概览}
投行/券商 27 篇；战略咨询 3 篇；智库/国际机构 9 篇。\par
投行覆盖：MS 9篇 / DB 5篇 / Citi 4篇 / GS 3篇 / JPM 2篇 / UBS 2篇 / BofA 1篇 / JEF 1篇\par
\section{Disclaimer}
{\scriptsize\color{refgray}
This community is a paid learning and information-sharing community created for general educational purposes only. The membership fee is charged solely for access to community discussions, educational content, organization, and operational costs. It is not an advisory fee, management fee, performance fee, brokerage fee, or compensation for personalized financial, investment, legal, tax, or professional advice.\par
I participate and share information solely as an individual and for educational discussion among community members. I am not acting as a financial adviser, investment adviser, broker, dealer, portfolio manager, legal adviser, tax adviser, accountant, fiduciary, or any other licensed professional adviser. Nothing shared in this community constitutes, and should not be interpreted as, financial advice, investment advice, legal advice, tax advice, a recommendation, solicitation, offer, endorsement, instruction, or guarantee to buy, sell, hold, short, trade, or invest in any security, cryptocurrency, fund, derivative, strategy, or other financial product.\par
All content shared in this community, including messages, discussions, links, files, charts, examples, market commentary, watchlists, AI-generated or AI-polished summaries, and third-party materials, is general, impersonal, and educational in nature. It does not take into account any individual's financial situation, investment objectives, risk tolerance, investment horizon, liquidity needs, tax status, legal restrictions, or suitability requirements. No content should be relied upon as suitable for any specific person, account, portfolio, or financial decision.\par
Information shared in this community may be incomplete, inaccurate, outdated, speculative, AI-generated, AI-edited, or based on third-party internet sources that have not been independently verified. I make no representation or warranty as to the accuracy, completeness, timeliness, reliability, or suitability of any information shared.\par
Financial markets involve substantial risk, including the possible loss of principal. Past performance, historical data, hypothetical examples, simulated results, backtests, personal opinions, or educational examples do not guarantee future results. Each member is solely responsible for conducting their own research, exercising independent judgment, and making their own decisions. Before making any financial, investment, legal, or tax decision, members should consult qualified and licensed professionals.\par
Participation in this community, payment of any membership fee, or communication with me or other members does not create any adviser-client, fiduciary, professional, contractual, agency, partnership, employment, or other advisory relationship. By joining or participating in this community, you acknowledge and agree that you are solely responsible for your own decisions, actions, transactions, profits, losses, risks, and consequences, and that I shall not be liable for any direct, indirect, incidental, consequential, financial, legal, tax, or other loss or damage arising from or related to any content, discussion, information, or material shared in this community.\par
}
\newpage
\begin{center}
{\Large 更多详情报告kcdesk.com}\par\vspace{0.8cm}
\includegraphics[width=0.58\linewidth]{\detokenize{../../prompts/zsxq_img.jpg}}
\end{center}
\end{document}
